\section{A certificate requirement is a different question from a compliant submission}
\label{sec:typed-decision-assurance}

Consider an XML suspicious transaction report submitted without the required
electronic certificate. Two correct answers are possible: the certificate
requirement applies, and the submission does not satisfy that requirement.
They disagree as Boolean values because they answer different questions. A
system that treats them as rival interpretations can reject a useful rule,
or select the wrong rule merely because it agrees with a test labelled
``compliance.'' Before comparing formulas, the investigation must identify what
the result is about.

SFC circular 26EC2 makes this distinction concrete. It lists three ways to
submit an STR through STREAMS~2: XML, upload of the prescribed PDF form, and
completion of the web form. Footnote~3 attaches the electronic-certificate
requirement to the first two methods. Footnote~2 separately requires reports
submitted through other channels after implementation to be resubmitted through
STREAMS~2. The transitional arrangements provide direct JFIU contact for urgent
reports during the blackout. Contact during the blackout, a compliant submission
after launch, and a duty to resubmit an earlier report are consequently separate
decisions \citep{sfcstr2026}.

For a deliberately limited example, suppose a licensed firm's original report
was submitted after the operational launch and its method is known to be exactly
one of the three listed methods. Let $x$, $p$ and $w$ denote selection of XML,
PDF and web form respectively. Let $s$ mean that the original event used
STREAMS~2, and $c$ mean that the certificate condition was fulfilled in the
accepted technical manner. These are facts about one event. Establishing the
firm's regulatory classification or the technical sufficiency of a certificate
requires evidence outside these Boolean formulas. Under those assumptions, the
certificate-applicability predicate $A$, the selected submission-control
predicate $C$, and the original-event resubmission trigger $R$ are
\begin{align}
 A &= x \lor p,\label{eq:str-certificate-applicability}\\
 C &= s \land \bigl(w \lor (x\land c) \lor (p\land c)\bigr),\label{eq:str-selected-submission}\\
 R &= \neg s.\label{eq:str-resubmission-trigger}
\end{align}
The scope conditions remain separate from these expressions. In particular,
$R$ is meaningful here only for an STR actually submitted to the JFIU after the
selected implementation boundary; it is not a duty imposed on every object
whose record lacks a STREAMS~2 flag.

For the XML example, $x=\mathsf{true}$, $s=\mathsf{true}$ and
$c=\mathsf{false}$ give $A=\mathsf{true}$, $C=\mathsf{false}$ and
$R=\mathsf{false}$. The failure to supply the certificate does not establish
that the report went through another channel. For a completed web-form
submission, $w=s=\mathsf{true}$ gives $A=\mathsf{false}$ and
$C=\mathsf{true}$ even if certificate possession is unknown. An unknown method
must not be represented by three false flags: the adapter preserves the method
as unknown. The evaluator then cannot conclude that the certificate condition
does not apply. Conflicting facts remain conflicts under the previously defined
semantics.

These formulas implement only the named conditions. The circular also discusses
XML schema conformity, technical testing, registration and transitional contact.
The combined output is a vector of the selected decisions, with their unresolved
dependencies, rather than an approval of the firm's reporting arrangements.
This preserves the practical benefit of an executable rule while preventing a
partial rule from inheriting the meaning of whole-document compliance.

\subsection{Checking a smaller question without losing the rest of the circular}

A typed control records its question, actor, unit of assessment, time basis,
result kind, true meaning and false meaning, together with source quotations.
For example, certificate applicability concerns one selected submission method,
whereas a report-level resubmission condition may quantify over several events.
Two rules are eligible for a direct Boolean comparison only after these
commitments have been aligned explicitly. Similar variable names or a common
result type are insufficient.

This distinction also changes source-fidelity checking. The earlier STR
investigation retained seventeen readings and forty-three claims, including
context. Checking every claim against every partial reading creates 731 pairs,
many of which compare different decisions. Two separately requested readers now
assign each claim and reading to the proposed controls. A scope definition may
qualify several controls. A reading that combines questions remains unresolved;
it is not silently narrowed. If the readers disagree, the affected comparison
remains required. A control with relevant claims but no candidate is reported
as missing. Claims originally labelled context are reconsidered because their
relevance can change with the question.

The assignments are themselves interpretations. Agreement does not prove that
a qualification is irrelevant, especially when both readers use the same model
family. Every excluded pair therefore retains both assignment reasons, and the
full source remains available to each fidelity request. A check can challenge
the assignment as well as the formula. This reduces avoidable work without
equating a shorter comparison list with completeness.

Each required pair is checked in a bounded batch. Completed responses, failed
attempts and pending identities are stored separately. If a batch cannot be
validated, the controller splits it and executes smaller requests within its
remaining allowance. It does not discard a difficult pair to obtain a complete
status. An interruption consumes its reservation; resumption reuses the exact
retained response only when its source, request and method bindings still match.
Thus the final result can distinguish a completed check that found an omission
from a check that never ran. Neither should be reported as source support.

\subsection{Repairing notation while preserving the proposed meaning}

Some retained model responses put an English explanation in a field intended
for a formal expression. Rejecting that response loses a potentially useful
interpretation, but asking a model simply to ``make it compile'' can change its
meaning. The repair must have a narrower purpose. It retains the entire original
reading, including factual definitions, classifications, assumptions, citations
and questions, and proposes only replacement scope and result expressions.

The expression compiler checks types and grammar. It also rejects an undeclared
fact rather than supplying a convenient default. The repaired expression is
then rendered back into controlled English. A separate context compares that
rendering with the original proposal and the full source in both directions:
what condition was lost, and what new condition was added? The correspondence
judgment and any uncertainty are retained beside the original. Even a successful
correspondence check supports a conditional claim about this translation. It
does not establish that the original English interpretation was legally correct.

This is especially useful for an ambiguous e-IP response. One proposed scope
described a test of whether a submission was made through e-IP, while its formal
body treated out-of-scope cases as automatically true. Moving the condition from
the result to the scope would change a true result into an out-of-scope result
for some inputs. That is a semantic decision, not cosmetic syntax repair. The
record must expose the difference before any corrected expression is promoted.

\subsection{New source evidence can settle a particular dispute}

Circular 25EC71 announced the equity-option margin exemption from 4 January
2026 while saying that the corresponding Code amendment would be gazetted in
due course. The retained alternatives disagreed about whether a transaction
therefore needed an additional gazettal-status fact. The January 2026 Code
provides more specific evidence: Schedule~10, Part~III, paragraph~7(e) lists the
non-centrally cleared single-stock, equity-basket and equity-index options, and
footnote~9 states the extension from 4 January 2026 until further notice.
The current provision supports the stated start date. It does not make a
separate transaction-level gazettal Boolean an expressed condition of that
provision. Historical publication timing and the possibility of a later notice
remain distinct questions \citep{sfcmargin2025,sfccode2026}.

The e-IP example illustrates a different repair. Circulars 24EC33 and 24EC50
both enumerate the investment-product categories administered by IPD. The later
circular expressly extends the parallel period and gives the mandatory channel
date. A rule can therefore expose the individual listed categories instead of
asking for an unexplained ``IPD product'' Boolean. A new product may still need
classification, but the missing judgment is now visible. Optional payment of
annual fees must remain distinguishable from an application or submission.
Likewise, a sourced extension of the parallel period justifies a priority
relationship for that particular date question; it does not make every sentence
of the later circular superior to every sentence of the earlier one
\citep{sfceiplaunch2024,sfceipadoption2024}.

An authority relationship consequently records the provisions it connects, the
question it affects, its applicable dates and the exact supporting passage.
A missing interval or a cycle leaves priority unresolved. The same discipline
applies to bilingual material. The retained English and traditional-Chinese
versions of 26EC2 can be paired at the certificate footnote, resubmission
footnote and operational-launch provision. Those pairings permit comparison;
matching reference numbers and dates do not prove equal meaning. Both texts and
the proposed correspondence remain available for later challenge.

The continuing resubmission duty adds a final temporal distinction. The selected
source does not supply a completion deadline. The event engine can record an
opening triggered by an original submission through another channel and a later
performance through STREAMS~2. Without a deadline, absence of performance leaves
the duty open; the engine must not invent a timed breach. If the dependent
authority changes, its current assurance is invalidated while the original
opening and historical replay remain bound to their original rule version.
Changing the law is not permission to rewrite the event history.

\subsection{What the executed demonstration establishes}

The development demonstration compiles the three STR decisions as separate
Java outputs and evaluates twelve explicitly described scenarios, including
missing methods, missing certificates and an incorrect original channel. A
second rule exposes the eight e-IP categories and its scope boundaries. Python,
ordinary generated Java and cvc5 compare the resulting forty-eight decisions;
the STR cases also run through the Catala backend. The duty example executes
both an open obligation and recorded performance, and checks that changed
authority requires reinvestigation. These are conditional execution checks on
authored cases. They are not forty-eight independent legal interpretations.

The PDF review has a similarly precise scope. It locates footnote attachments
and the four return or yield columns in the retained advertising examples.
Exact relocation of a visually inspected running page label can be resolved
without altering legal wording. Changed quantities, units, negations, exceptions
or column labels remain discrepancies. An unresolved OCR difference is not
converted into agreement merely because another parser repeats it. The
development annotations preserve the page, coordinates, source hash and image
hash so that a reviewer can inspect the particular proposed resolution.

These mechanisms make the investigation more specific and more useful. They
also expose what the present evidence does not measure: independent legal
accuracy across unseen source families, error correlation between different
model families, and actual reviewer time saved. Those questions require suitable
reference judgments and observations. They must not be answered by counting
compiled rules, matching outputs, resolved page labels or repeated model votes.

The gifts investigation illustrates that distinction in the executed record.
It retained all fifty-three claims and eight readings, giving 424 possible
claim--reading pairs. The two routing judgments proposed excluding 203 pairs
because they concerned different questions. All 221 remaining checks completed:
the model judged ninety entailed and 131 not established, and raised ninety-two
additional concerns. Several questions still required external definitions or
authority. These counts establish completed work and the character of the
findings; they do not measure the model's legal accuracy. In particular, an
unestablished pair may identify a missing condition, an unencoded reading, or
insufficient evidence rather than a legally false proposition.

The STR run exposed a different problem: a reader used local proposal names
instead of inventory identities and subsequently omitted two claims. Both raw
responses were retained. The repair requests smaller portions of the inventory
while keeping the complete circular and control definitions in every request.
Only an exactly accounted union can support routing. A resource ceiling may
still leave checks pending, but it cannot make those checks disappear from the
reported universe. This is the operational distinction between an investigation
that has finished its allowed work and an interpretation that has been settled.
